\begin{table}[htbp]
\centering
\caption{Descriptive statistics for the main analysis sample}
\label{tab:sample_description}
\small
\begin{threeparttable}
\begin{adjustbox}{max width=\textwidth}
\begin{tabular}{lrrrrrrrr}
\toprule
Variable & $N$ & Mean & Variance & Min & Median & Max & Urban mean & Rural mean \\
\midrule
Food share & 23,971 & 0.378 & 0.027 & 0.000 & 0.364 & 0.997 & 0.341 & 0.421 \\
Urban treatment & 23,971 & 0.531 & 0.249 & 0.000 & 1.000 & 1.000 & 1.000 & 0.000 \\
Log total expenditure & 23,971 & 13.435 & 0.522 & 9.589 & 13.502 & 16.249 & 13.590 & 13.259 \\
Age & 23,971 & 50.540 & 228.213 & 16.000 & 50.000 & 99.000 & 48.784 & 52.525 \\
Household size & 23,971 & 3.694 & 3.165 & 1.000 & 4.000 & 17.000 & 3.717 & 3.667 \\
Male reference person & 23,971 & 0.860 & 0.120 & 0.000 & 1.000 & 1.000 & 0.847 & 0.876 \\
Propensity score & 23,971 & 0.531 & 0.019 & 0.060 & 0.537 & 0.943 & 0.567 & 0.489 \\
\bottomrule
\end{tabular}
\end{adjustbox}
\begin{tablenotes}[flushleft]
\footnotesize
\item The table reports the main analysis variables after the common treatment-effect sample restrictions are imposed. Urban treatment is defined as \(1(\texttt{town} \ge 4)\).
\end{tablenotes}
\end{threeparttable}
\end{table}
